% !TeX program = xelatex
\documentclass[12pt,a4paper]{article}
\usepackage{amsmath,amssymb}
\usepackage{geometry}
\usepackage{booktabs}
\usepackage{graphicx}
\usepackage{xcolor}
\usepackage[UTF8]{ctex}
\geometry{margin=2.5cm}
\title{Q2 Logistic Model: Symmetric Antagonistic $\alpha$, $\beta$}
\author{}
\date{}

\begin{document}
\maketitle

\section{Model Formulation}

为进一步提高基础 Logistic 模型对实际恢复过程的描述能力，
本文在 Q1 模型基础上引入两个具有拮抗作用的调节参数：
促进系数 $\alpha$ 与抑制系数 $\beta$。

其中，$\alpha$ 表征外部促进因素对系统潜在恢复能力的增强作用，
通过提高有效环境容量 $K_{\mathrm{eff}}$ 改变系统长期稳定水平；
$\beta$ 表征内部限制因素对恢复过程的约束作用，
通过降低有效增长率 $r$ 调节系统接近稳定状态的速度。

二者作用方向相反，但共同决定系统动态演化过程，
因此形成具有促进—抑制拮抗机制的扩展 Logistic 模型：

\begin{equation}\label{eq:ode}
\frac{dx}{dt} = \frac{\gamma}{1+\beta} \cdot x \cdot \left(1 - \frac{x}{K(1+\alpha)}\right)
\end{equation}

\begin{itemize}
  \item $\alpha\ge0$ 表示外部促进作用，其通过提高系统潜在承载能力扩大有效环境容量：  $\;\rightarrow\; K_{\text{eff}} = K(1+\alpha)$
  \item $\beta\ge0$ 表示内部抑制作用，其通过降低系统有效增长率影响恢复速度： $\;\rightarrow\; r = \gamma/(1+\beta)$
\end{itemize}

两个调节参数均采用 $(1+c)^{\pm1}$ 的形式作用于模型参数，
从数学结构上保持参数调整的对称性。

其中，$\alpha$ 主要影响系统长期稳定水平，
而 $\beta$ 主要影响系统达到稳定状态的速度。

当 $\alpha=\beta=0$ 时，扩展模型退化为 Q1 中的基础 Logistic 模型。拮抗方向相反：
$\alpha$ 增大饱和水平，$\beta$ 降低趋近速度。当 $\alpha = \beta = 0$ 时退化为 Q1。

\section{Standard Logistic Form}

经过参数变换后，扩展模型仍保持 Logistic 模型的标准形式，
但其有效环境容量和有效增长率均由调节参数动态决定。

因此，该模型既继承了 Logistic 模型的稳定收敛特性，
又能够描述外部促进因素和内部限制因素共同作用下的恢复过程。

\begin{equation}\label{eq:logistic}
\frac{dx}{dt} = r \cdot x \cdot \left(1 - \frac{x}{K_{\text{eff}}}\right)
\end{equation}
\begin{equation}
r = \frac{\gamma}{1+\beta}, \qquad K_{\text{eff}} = K(1+\alpha)
\end{equation}

解析解：
\begin{equation}\label{eq:solution}
x(t) = \frac{K_{\text{eff}}}{1 + \left(\dfrac{K_{\text{eff}}}{x_0} - 1\right) e^{-r t}}
\end{equation}

\section{Parameter Identification}

由于扩展模型最终可以转化为标准 Logistic 形式，
因此本文首先对有效增长率 $r$ 和有效环境容量 $K_{eff}$ 进行参数识别。

完成参数拟合后，再根据参数映射关系反推出促进系数 $\alpha$
和抑制系数 $\beta$：
\begin{equation}
\beta = \frac{\gamma}{r} - 1, \qquad \alpha = \frac{K_{\text{eff}}}{K} - 1
\end{equation}

根据模型设计要求，当拟合结果满足
$r\le\gamma$ 且 $K_{eff}\ge K$ 时，
可进一步保证 $\beta\ge0$ 和 $\alpha\ge0$。

该约束条件符合实际恢复过程中的物理意义，
即促进因素不会降低系统容量，而抑制因素不会提高增长速率。

\section{Parameters}

\begin{center}
\begin{tabular}{llrl}
\toprule
类别 & 符号 & 拟合值 & 单位 \\
\midrule
固定          & $K$                & 8011        & tons/day \\
固定          & $\gamma$           & 0.022       & 1/day \\
固定          & $x_0$              & 6314        & tons/day \\
\midrule
\textbf{拟合} & $r$                & $\mathbf{0.008485} \pm 0.000655$ & 1/day \\
\textbf{拟合} & $K_{\text{eff}}$   & $\mathbf{8142.22} \pm 72.77$ & tons/day \\
\midrule
导出          & $\alpha$           & $\mathbf{0.016379}$ & --- \\
导出          & $\beta$            & $\mathbf{1.592906}$ & --- \\
\bottomrule
\end{tabular}
\end{center}

\textbf{参数物理意义分析}：
\begin{itemize}
  \item $\alpha=0.0164>0$ 表明系统存在正向促进作用，
使有效环境容量提高约 $1.64\%$，
即 $K_{eff}$ 略高于基础模型中的理论容量 $K$。
  \item $\beta=1.5929>0$ 表明系统受到明显抑制作用，
有效增长率降低至
$r=\gamma/(1+\beta)=0.0085$，
约为原始增长率的 $38.6\%$。

该结果说明实际恢复过程中的增长速度受到较强限制。
  \item $r/\gamma = 1/(1+\beta) = 0.386$：综合分析可知，
增长率修正幅度明显大于容量修正幅度，
说明在当前恢复阶段，内部限制因素对系统演化过程起主要作用。
\end{itemize}

\section{Evaluation}

为了验证扩展模型相对于基础 Logistic 模型的改进效果，
本文将 Q1 基础模型与引入 $\alpha,\beta$ 参数后的扩展模型进行对比。

通过比较两类模型的拟合指标，
分析调节参数对于提高预测精度和降低系统误差的有效性。

\begin{center}
\begin{tabular}{lcc}
\toprule
Metric & Q1 (no $\alpha,\beta$) & Q2 (with $\alpha,\beta$) \\
\midrule
$R^2$  & 0.5319 & \textbf{0.9915} \\
MAE    & 315.6  & \textbf{38.3} \\
MAPE   & 4.42\% & \textbf{0.53\%} \\
\bottomrule
\end{tabular}
\end{center}

从预测结果可以看出，引入 $\alpha$ 和 $\beta$ 后，
模型预测值与实际观测值之间的偏差明显降低。

相比 Q1 中持续存在的系统性高估现象，
扩展模型在整个时间范围内均保持较小残差，
说明调节参数能够有效修正基础 Logistic 模型中增长率设定偏高的问题。

\section{Predictions}

\begin{center}
\begin{tabular}{crrrr}
\toprule
$t$ (days) & Actual & Predicted & Residual & Rel.\ Error \\
\midrule
0   & 6314 & 6314.0 & $-0.0$    & $-0.00\%$ \\
30  & 6542 & 6649.5 & $-107.5$  & $-1.64\%$ \\
60  & 6875 & 6935.3 & $-60.3$   & $-0.88\%$ \\
90  & 7173 & 7174.2 & $-1.2$    & $-0.02\%$ \\
120 & 7368 & 7371.2 & $-3.2$    & $-0.04\%$ \\
150 & 7591 & 7531.5 & $+59.5$   & $+0.78\%$ \\
180 & 7724 & 7660.6 & $+63.4$   & $+0.82\%$ \\
270 & 7896 & 7910.5 & $-14.5$   & $-0.18\%$ \\
365 & 8002 & 8037.1 & $-35.1$   & $-0.44\%$ \\
\bottomrule
\end{tabular}
\end{center}

\section{Key Findings}

\begin{itemize}

\item
\textbf{促进与抑制机制同时存在：}
拟合结果表明 $\alpha=0.0164>0$，
$\beta=1.5929>0$，
说明恢复系统同时受到外部促进作用和内部限制作用影响。

\item
\textbf{参数作用具有明显分工：}
促进系数 $\alpha$ 主要影响系统长期稳定容量，
使有效容量提高至 $K_{eff}=8142.22$ tons/day；
抑制系数 $\beta$ 主要影响恢复速度，
使有效增长率降低至原始增长率的 $38.6\%$。

\item
\textbf{模型预测能力显著提升：}
引入调节参数后，
模型决定系数由 Q1 的 $0.5319$ 提升至 $0.9915$，
MAPE 从 $4.42\%$ 降低至 $0.53\%$，
说明扩展 Logistic 模型能够更加准确地描述实际恢复过程。

\item
\textbf{模型具有良好的稳定性：}
当 $\alpha,\beta\ge0$ 时，
有效增长率保持正值且系统仍满足 Logistic 收敛性质，
不会发生失稳现象。

\item
\textbf{后续研究方向：}
虽然当前模型已经显著提高预测精度，
但 $\alpha$ 与 $\beta$ 仍被假设为常数。
未来可进一步考虑时间变化参数，
以描述不同恢复阶段中促进和抑制作用的动态变化。

\end{itemize}

\end{document}
